% Options for packages loaded elsewhere
\PassOptionsToPackage{unicode}{hyperref}
\PassOptionsToPackage{hyphens}{url}
\PassOptionsToPackage{dvipsnames,svgnames,x11names}{xcolor}
\documentclass[
  10pt,
]{ctexart}
\usepackage{xcolor}
\usepackage[margin=2cm]{geometry}
\usepackage{amsmath,amssymb}
\setcounter{secnumdepth}{-\maxdimen} % remove section numbering
\usepackage{iftex}
\ifPDFTeX
  \usepackage[T1]{fontenc}
  \usepackage[utf8]{inputenc}
  \usepackage{textcomp} % provide euro and other symbols
\else % if luatex or xetex
  \usepackage{unicode-math} % this also loads fontspec
  \defaultfontfeatures{Scale=MatchLowercase}
  \defaultfontfeatures[\rmfamily]{Ligatures=TeX,Scale=1}
\fi
\usepackage{lmodern}
\ifPDFTeX\else
  % xetex/luatex font selection
  \setmainfont[]{DejaVu Sans}
\fi
% Use upquote if available, for straight quotes in verbatim environments
\IfFileExists{upquote.sty}{\usepackage{upquote}}{}
\IfFileExists{microtype.sty}{% use microtype if available
  \usepackage[]{microtype}
  \UseMicrotypeSet[protrusion]{basicmath} % disable protrusion for tt fonts
}{}
\makeatletter
\@ifundefined{KOMAClassName}{% if non-KOMA class
  \IfFileExists{parskip.sty}{%
    \usepackage{parskip}
  }{% else
    \setlength{\parindent}{0pt}
    \setlength{\parskip}{6pt plus 2pt minus 1pt}}
}{% if KOMA class
  \KOMAoptions{parskip=half}}
\makeatother
\usepackage{longtable,booktabs,array}
\usepackage{caption}
\captionsetup[table]{skip=6pt}
\newcounter{none} % for unnumbered tables
\usepackage{calc} % for calculating minipage widths
% Correct order of tables after \paragraph or \subparagraph
\usepackage{etoolbox}
\makeatletter
\patchcmd\longtable{\par}{\if@noskipsec\mbox{}\fi\par}{}{}
\makeatother
% Allow footnotes in longtable head/foot
\IfFileExists{footnotehyper.sty}{\usepackage{footnotehyper}}{\usepackage{footnote}}
\makesavenoteenv{longtable}
\setlength{\emergencystretch}{3em} % prevent overfull lines
\providecommand{\tightlist}{%
  \setlength{\itemsep}{0pt}\setlength{\parskip}{0pt}}
\usepackage{bookmark}
\IfFileExists{xurl.sty}{\usepackage{xurl}}{} % add URL line breaks if available
\urlstyle{same}
% fallback for those not using the hyperref driver hyperxmp:
\makeatletter
\@ifundefined{xmpquote}{\newcommand{\xmpquote}[1]{#1}}{}
\makeatother
\hypersetup{
  colorlinks=true,
  linkcolor={Maroon},
  filecolor={Maroon},
  citecolor={Blue},
  urlcolor={Blue},
  pdfcreator={LaTeX via pandoc}}

\author{}
\date{}

\usepackage{pdflscape,fvextra}
\setlength{\tabcolsep}{3pt}
\begin{document}

\section{LUAD Top-10
对疾病签名阈值的敏感性（v1）}\label{luad-top-10-ux5bf9ux75beux75c5ux7b7eux540dux9608ux503cux7684ux654fux611fux6027v1}

2026-10-01。冻结排名用 RRF（k=60）融合两项：所有匹配 landmark 基因上的负
Spearman
反转（与阈值无关），以及上/下调基因集连接性（依赖阈值）。这里只重建连接性基因集：阈值网格为
FDR ∈ \{0.01, 0.05, 0.10\} × \textbar log2FC\textbar{} ∈ \{0.58, 1.0,
1.5\}；另加每个方向按 \textbar log2FC\textbar{} 取前 N 个（N =
25、50、100，FDR \textless{} 0.05）。默认设定（FDR
0.05、\textbar log2FC\textbar{} 1）精确复现冻结
Top-10。冻结排名与证据不因本分析改变。脚本：\texttt{python\ -m\ evals.\allowbreak{}luad\_\allowbreak{}threshold\_\allowbreak{}sensitivity}；结果：\texttt{benchmark/\allowbreak{}results/\allowbreak{}luad\_\allowbreak{}threshold\_\allowbreak{}sensitivity\_\allowbreak{}v1.\allowbreak{}json}。糖皮质激素按名称规则判定，规则与逐个标签都写在
JSON 中；其他药物的机制取自 Broad Repurposing Hub 2025-08-18 注释。

{\def\LTcaptype{none} % do not increment counter
\begin{longtable}[]{@{}>{\raggedright\arraybackslash}p{0.1800\textwidth}>{\raggedright\arraybackslash}p{0.1800\textwidth}>{\raggedright\arraybackslash}p{0.1800\textwidth}>{\raggedright\arraybackslash}p{0.1800\textwidth}>{\raggedright\arraybackslash}p{0.1800\textwidth}@{}}
\toprule\noalign{}
设定 & 上/下调 landmark 基因 & 与冻结 Top-10 重合 & Top-10 中糖皮质激素
& 参考药平均百分位 \\
\midrule\noalign{}
\endhead
\bottomrule\noalign{}
\endlastfoot
FDR\textless0.01, & log2FC & \textgreater=0.58 & 155/121 & 3/10 \\
FDR\textless0.01, & log2FC & \textgreater=1 & 57/50 & 10/10 \\
FDR\textless0.01, & log2FC & \textgreater=1.5 & 16/18 & 7/10 \\
FDR\textless0.05, & log2FC & \textgreater=0.58 & 155/121 & 3/10 \\
FDR\textless0.05, & log2FC & \textgreater=1 & 57/50 & 10/10 \\
FDR\textless0.05, & log2FC & \textgreater=1.5 & 16/18 & 7/10 \\
FDR\textless0.1, & log2FC & \textgreater=0.58 & 155/121 & 3/10 \\
FDR\textless0.1, & log2FC & \textgreater=1 & 57/50 & 10/10 \\
FDR\textless0.1, & log2FC & \textgreater=1.5 & 16/18 & 7/10 \\
top-25 per direction (FDR\textless0.05) & 25/25 & 8/10 & 9 & 0.386 \\
top-50 per direction (FDR\textless0.05) & 50/50 & 9/10 & 9 & 0.361 \\
top-100 per direction (FDR\textless0.05) & 100/100 & 4/10 & 4 & 0.402 \\
\end{longtable}
}

\textbf{结论}

\begin{itemize}
\tightlist
\item
  \textbf{FDR 阈值不起作用}：在 961 个 landmark
  基因中，\textbar log2FC\textbar{} ≥ 0.58 的基因全部已满足 FDR
  \textless{} 0.01，三档 FDR 给出完全相同的结果。
\item
  \textbf{\textbar log2FC\textbar{} 阈值决定 Top-10
  的面貌}：默认与更严的 1.5 下，Top-10 中有 9 个糖皮质激素；放宽到 0.58
  或每方向取前 100 个基因时，与冻结 Top-10 仅重合 3--4
  个，糖皮质激素降到 3--4 个。
\item
  放宽阈值（0.58）后的 Top-10：AZ-628（RAF
  inhibitor）、clocortolone-pivalate（糖皮质激素）、beclomethasone-dipropionate（糖皮质激素）、hydrocortisone（糖皮质激素）、importazole（Hub
  无注释）、GDC-0941（PI3K inhibitor）、BRD-K97972722（Hub
  无注释）、wortmannin（PI3K inhibitor）、NVP-BEZ235（mTOR inhibitor
  \textbar{} PI3K inhibitor）、BRD-A47816767（Hub 无注释）。新进入者以
  RAF、PI3K、PI3K/mTOR 抑制剂为主，对应 A549（KRAS 突变）中的
  RAF--PI3K--mTOR 信号轴。这在生物学上可解释，但同样只是转录组假设。
\item
  12 种设定下始终在 Top-10
  的只有：beclomethasone-dipropionate、clocortolone-pivalate、hydrocortisone。
\item
  参考药的平均百分位在各设定下为
  0.361--0.427，始终没有明显优于随机（0.5）。
\end{itemize}

\textbf{对答辩表述的影响}：``Top-10 中 9/10
为糖皮质激素''只在预设阈值（\textbar log2FC\textbar{} ≥
1）及更严的设定下成立，不是稳健结论。准确的说法是：糖皮质激素信号在严格签名下占主导，放宽签名后被
RAF/PI3K/mTOR 抑制剂部分取代；只有 3 个糖皮质激素在所有设定下都留在
Top-10。预设阈值在看排名之前就已固定，因此主结果不改；本分析作为稳健性披露。

\end{document}
